% Source truth: source/普通高考/2010/2010广东理.pdf, p.1, q17 histogram.
% Data: intervals (490,495], (495,500], (500,505], (505,510], (510,515]
% have frequency-density heights 0.03, 0.04, 0.07, 0.05, 0.01 (from the source graph).
% solid segments: five unfilled bar outlines; dashed segments: horizontal guides at
% 0.01, 0.03, 0.04, 0.05, 0.07; labels: fraction on the y-axis, weight/unit on x-axis.
\begin{tikzpicture}[baseline=(current bounding box.center)]
\begin{axis}[exam statistical histogram,
  width=9.2cm,
  height=5.6cm,
  xmin=487,
  xmax=518,
  ymin=0,
  ymax=0.08,
  scaled y ticks=false,
  axis lines=left,
  axis line style={-Stealth},
  tick align=outside,
  xtick={487,490,495,500,505,510,515},
  xticklabels={0,490,495,500,505,510,515},
  ytick={0.01,0.03,0.04,0.05,0.07},
  yticklabel style={/pgf/number format/fixed, /pgf/number format/precision=2},
  tick label style={font=\normalsize},
  label style={font=\normalsize},
  clip=false,
]
    \examstatxlabel{重量/克};
  % Source dashed guides terminate at the corresponding bar boundary.
  \addplot[black, dashed, line width=0.45pt, forget plot]
    coordinates {(487,0.01) (515,0.01)};
  \addplot[black, dashed, line width=0.45pt, forget plot]
    coordinates {(487,0.03) (495,0.03)};
  \addplot[black, dashed, line width=0.45pt, forget plot]
    coordinates {(487,0.04) (500,0.04)};
  \addplot[black, dashed, line width=0.45pt, forget plot]
    coordinates {(487,0.05) (505,0.05)};
  \addplot[black, dashed, line width=0.45pt, forget plot]
    coordinates {(487,0.07) (500,0.07)};
  % Frequency-density bars, drawn as unfilled outlines to preserve visible guides.
  \addplot[black, ybar interval, draw=black, fill=none, line width=0.55pt]
    coordinates {(490,0.03) (495,0.04) (500,0.07) (505,0.05) (510,0.01) (515,0)};
  % The source x-axis uses a short break immediately after the origin.
  \draw[white, line width=2.2pt] (axis cs:487.5,0) -- (axis cs:489.2,0);
  \examstatxbreak[black, line width=0.55pt]{488.4}{0};
  % Upright source-style fraction label above the y-axis.
  \examstatylabel{\(\dfrac{\text{频率}}{\text{组距}}\)};
\end{axis}
\end{tikzpicture}
